\chapter{What is Risk ?}
\label{chap:What is Risk}

%---------------------------- SECTION ----------------------------
\section{The Word We Use Without Always Understanding}

\paragraph{Risk is one of those words that appears constantly in professional life. It turns up in board reports, regulatory filings, contract negotiations, and casual conversations alike. We talk about risking a decision, about something being too risky, or about the risk of not acting at all. Yet for a word so frequently used, it is surprisingly rarely defined — and that lack of definition can cause real problems.}

\paragraph{In a professional context, and particularly in compliance and legal settings, the word risk carries specific and important meaning. Getting that meaning right from the outset is not merely an academic exercise. It shapes how your organisation identifies problems, makes decisions, allocates resources, and ultimately demonstrates to regulators and stakeholders that it is operating responsibly.}

\paragraph{So let us start at the very beginning.}

%---------------------------- SECTION ----------------------------
\section{Defining Risk}

\paragraph{At its simplest, risk can be understood as the possibility that something will happen that has an impact on your objectives.}

\paragraph{That definition is deliberately broad — and intentionally so. Risk is not always negative. In financial services, risk is often talked about in terms of exposure to potential loss, but risk also encompasses the upside — the chance that a decision pays off, that a new product succeeds, or that entering a new market generates growth.}

\paragraph{However, for the purposes of compliance and legal risk assessment, we are primarily — though not exclusively — concerned with the downside: the things that could go wrong, the obligations that could be breached, and the consequences that could follow.}

\paragraph{A more formal definition, drawn from the internationally recognised standard ISO 31000, describes risk as: }

\barquote{The effect of uncertainty on objectives.}

\paragraph{This is a deceptively simple but powerful framing. Let's unpack it:}

\begin{itemize}
  \bitem{Effect}{risk is about impact. Something happens, and it has consequences. Those consequences might be financial, legal, operational, reputational, or human.}
  \bitem{Uncertainty}{risk exists precisely because we cannot know the future with certainty. If an outcome were guaranteed, it would not be a risk — it would simply be a fact to be managed.}
  \bitem{Objectives}{this is the part that is most often overlooked. Risk is always relative to what you are trying to achieve. A risk that is significant for one organisation may be inconsequential for another, depending entirely on their respective goals, structures, and obligations.}
\end{itemize}

\paragraph{For compliance professionals, those objectives are often defined at least in part by external forces — regulations, legislation, industry standards, and the expectations of regulators and customers. Understanding risk in your context therefore requires first understanding what your organisation is obligated and expected to achieve.}

%---------------------------- SECTION ----------------------------
\section{Hazard, Risk, and Consequence — Knowing the Difference}

\paragraph{Before going any further, it is worth drawing a distinction between three terms that are sometimes used interchangeably but which mean quite different things. Confusing them is one of the most common mistakes made by those new to formal risk assessment.}

\subsection{Hazard}

\paragraph{A hazard is a source of potential harm. It is the thing — whether a situation, an action, a condition, or a failure — that has the potential to cause a negative outcome. A hazard exists regardless of whether anyone is exposed to it.}

\example{Example}{A poorly drafted contract clause is a hazard. It sits there, waiting.}

\subsection{Risk}

\paragraph{Risk is the likelihood that the hazard will actually lead to harm, combined with the severity of that harm if it does. Risk only becomes real when there is both a hazard and exposure to it.}

\example{Example}{If that poorly drafted contract clause governs a high-value transaction with an aggressive counterparty, the risk is significant. If it appears in a rarely used template that has never been executed, the risk is lower — though not necessarily zero.}

\subsection{Consequence}

\paragraph{Consequence is what actually happens when a risk materialises. It is the outcome — the financial penalty, the regulatory sanction, the reputational damage, the litigation, or the operational disruption.}

\example{Example}{The counterparty exploits the ambiguous clause and your organisation faces unexpected liability.}

\paragraph{Understanding these distinctions matters because effective risk assessment requires you to think about all three. Identifying hazards alone is not enough. You need to understand exposure, likelihood, and potential impact before you can make informed decisions about what to do next.}

%---------------------------- SECTION ----------------------------
\newpage
\section{Risk vs. Uncertainty — Knowing the Difference}

\paragraph{There is a further distinction worth introducing early, one that becomes increasingly relevant as risk assessments grow more sophisticated: the difference between risk and uncertainty.}

\begin{itemize}
  \bitem{Risk}{Risk refers to situations where the possible outcomes are known — or can be reasonably estimated — and where we can assign some probability to each. We may not know what will happen, but we have enough information to make an informed judgement.}
  \bitem{Uncertainty}{Uncertainty refers to situations where we cannot reliably identify the possible outcomes, let alone assign probabilities to them. These are the unknown unknowns — the scenarios that fall outside our existing mental models or experience.}
\end{itemize}

\paragraph{In practice, most real-world compliance challenges sit somewhere on a spectrum between pure risk and pure uncertainty. Regulatory change, geopolitical shifts, emerging technologies, and novel business models all introduce elements of uncertainty that resist simple probabilistic analysis.}

\paragraph{This is not a reason to abandon structured risk assessment. Rather, it is a reminder that risk assessment is a tool for improving decision-making under conditions of imperfect knowledge — not a mechanism for predicting the future with precision. We will return to this theme in depth in Chapter~\ref{chap:Risk Assessment Under Uncertainty}.}

%---------------------------- SECTION ----------------------------
\section{Why the Language of Risk Matters}

\paragraph{You might wonder whether these definitional distinctions are really worth dwelling on. In practice, does it matter whether someone says hazard when they mean risk?}

\paragraph{The short answer is: yes, it can.}

\paragraph{When organisations lack a shared vocabulary for risk, several problems tend to follow:}

\begin{itemize}
  \bitem{Miscommunication}{Miscommunication between teams leads to the same issue being assessed differently in different parts of the business.}
  \bitem{Inconsistent reporting}{Inconsistent reporting makes it difficult for boards and senior leaders to get an accurate picture of the organisation's risk profile.}
  \bitem{Gaps in accountability}{Gaps in accountability emerge when it is unclear who owns a given risk versus who is simply aware of it.}
  \bitem{Regulatory submissions}{Regulatory submissions that are internally inconsistent can attract scrutiny and undermine credibility with oversight bodies.}
\end{itemize}

\paragraph{Establishing a clear, shared language for risk — one that is consistently used across your compliance framework, your documentation, and your communications — is a small investment that pays significant dividends over time. Many organisations achieve this through a risk taxonomy: a structured list of defined terms and risk categories used consistently across the business. We will look at how to build one in Chapter~\ref{chap:Step 1 — Identifying Hazards and Risk Sources}.}

 %---------------------------- SECTION ----------------------------
\newpage
\section{Risk Is Everywhere — But That Is Not a Reason for Paralysis}

\paragraph{It is easy, particularly in compliance-heavy environments, to become overwhelmed by the sheer volume and variety of risks an organisation faces. Every process, every relationship, every contract, every system carries some degree of risk. Taken to its logical extreme, risk awareness can feel less like a management tool and more like a source of paralysis.}

\paragraph{The purpose of risk assessment is not to eliminate risk — that is neither possible nor desirable. Organisations exist, in part, precisely to take on and manage risk in pursuit of value. A business that takes no risks generates no returns. A compliance programme that treats every conceivable risk as equally urgent will quickly exhaust itself and lose credibility with the business it is meant to support.}

\paragraph{The purpose of risk assessment is to ensure that risks are:}

\begin{itemize}
  \bitem{Understood reporting}{surfaced and acknowledged rather than ignored or assumed away.}
  \bitem{Known}{analysed with sufficient rigour to inform good decision-making.}
  \bitem{Proportionate}{responded to in a way that reflects their actual significance.}
  \bitem{Owned}{assigned to someone with the authority and capability to manage them.}
\end{itemize}

\paragraph{This book is about building the knowledge and practical capability to do exactly that. And it starts here — with understanding what risk actually is.}

 %---------------------------- CHAPTER SUMMARY ----------------------------
\section{Chapter Summary}
\paragraph{Before moving on, let us anchor the key ideas from this chapter:}

\paragraph{Key takeaways:}

\begin{itemize}
  \bitem{Attitude}{Risk is not inherently negative, but in compliance contexts we are primarily focused on downside risk.}
  \bitem{Terminology}{Hazard, risk, and consequence are distinct concepts — confusing them undermines effective assessment.}
  \bitem{Objectives}{Risk always exists relative to objectives — define those first.}
  \bitem{Results}{The goal is not to eliminate risk, but to know it, understand it, and manage it proportionately.}
  \bitem{Taxonomy}{A shared language for risk across your organisation is foundational to everything that follows. In this chapter we defined:
  \begin{itemize}
          \bitem{Risk}{The effect of uncertainty on objectives.}
          \bitem{Hazard}{A source of potential harm.}
          \bitem{Consequence}{The outcome when a risk materialises.}
          \bitem{Uncertainty}{Situations where outcomes and probabilities cannot be reliably estimated.}
        \end{itemize}
    }
\end{itemize}